\chapter{Gesättigte Teilmengen}
\label{b11guide:ch:saturation}

\section{Klassen vollständig aufnehmen}

\begin{definition}[Gesättigte Teilmenge]
Seien $E\in\EqSet(A)$ und $B\subseteq A$.
$B$ heißt \emph{$E$-gesättigt}, wenn
\[
 \forall x\in B\;\forall y\in A\;
          (x\sim_E y\Rightarrow y\in B).
\]
Die Familie aller $E$-gesättigten Teilmengen von $A$ wird mit
$\SatSets_E(A)$ bezeichnet.
\end{definition}

\begin{definition}[Sättigung und gesättigter Kern]
\label{b11guide:def:saturation}
Für $E\in\EqSet(A)$ und $B\subseteq A$ seien
\[
\begin{aligned}
 \Sat_E(B)&=\bigcup_{x\in B}[x]_E
           =\{y\in A\mid \exists x\in B\;(x\sim_E y)\},\\
 \Core_E(B)&=\bigcup\{C\in A/E\mid C\subseteq B\}.
\end{aligned}
\]
$\Sat_E(B)$ heißt die \emph{Sättigung} von $B$,
$\Core_E(B)$ sein \emph{gesättigter Kern}.
Diese Schreibweisen bezeichnen die angegebenen Mengen;
es wird keine Theorie von Abbildungen oder topologischen Räumen
vorausgesetzt.
\end{definition}

\begin{theorem}[Charakterisierungen gesättigter Mengen]
\label{b11guide:thm:saturated}
Für $E\in\EqSet(A)$ und $B\subseteq A$ sind äquivalent:
\[
\begin{aligned}
 &B\text{ ist }E\text{-gesättigt},\\
 &B=\Sat_E(B),\qquad B=\Core_E(B),\\
 &\exists\mathcal D\subseteq A/E\;(B=\bigcup\mathcal D),\\
 &\forall C\in A/E\;(C\subseteq B\ \lor\ C\cap B=\varnothing).
\end{aligned}
\]
Gesättigtheit ist somit genau die Eigenschaft, ganze Klassen
aufzunehmen oder ganz auszulassen.
\end{theorem}

\begin{theorem}[Kleinste gesättigte Erweiterung und größter gesättigter Teil]
\label{b11guide:thm:bounds}
Für $E\in\EqSet(A)$ und $B\subseteq A$ sind
$\Sat_E(B)$ und $\Core_E(B)$ gesättigt, und
\[
 \Core_E(B)\subseteq B\subseteq\Sat_E(B).
\]
Für jede gesättigte Menge $D\subseteq A$ gilt
\[
 B\subseteq D\ \Longrightarrow\ \Sat_E(B)\subseteq D,
 \qquad
 D\subseteq B\ \Longrightarrow\ D\subseteq\Core_E(B).
\]
\end{theorem}

\section{Rechengesetze der Sättigung}

\begin{theorem}[Wiederholung, Inklusion und Randfälle]
Für $E\in\EqSet(A)$ und $B\subseteq D\subseteq A$ gilt
\[
\begin{aligned}
 \Sat_E(B)&\subseteq\Sat_E(D),&
 \Core_E(B)&\subseteq\Core_E(D),\\
 \Sat_E(\Sat_E(B))&=\Sat_E(B),&
 \Core_E(\Core_E(B))&=\Core_E(B),\\
 \Sat_E(\varnothing)&=\varnothing,&
 \Core_E(\varnothing)&=\varnothing,\\
 \Sat_E(A)&=A,& \Core_E(A)&=A.
\end{aligned}
\]
Für $x\in A$ gilt $\Sat_E(\{x\})=[x]_E$.
\end{theorem}

\begin{theorem}[Komplemente und beliebige Familien]
\label{b11guide:thm:satlaws}
Seien $E\in\EqSet(A)$, $B\subseteq A$ und
$(B_i)_{i\in I}$ eine Familie von Teilmengen von $A$. Dann gilt
\[
\begin{aligned}
 \Core_E(B)&=A\setminus\Sat_E(A\setminus B),\\
 \Sat_E(B)&=A\setminus\Core_E(A\setminus B),\\
 \Sat_E\left(\bigcup_{i\in I}B_i\right)
       &=\bigcup_{i\in I}\Sat_E(B_i),\\
 \Core_E\left(\bigcap_{i\in I}B_i\right)
       &=\bigcap_{i\in I}\Core_E(B_i),\\
 \Sat_E\left(\bigcap_{i\in I}B_i\right)
       &\subseteq\bigcap_{i\in I}\Sat_E(B_i),\\
 \bigcup_{i\in I}\Core_E(B_i)
       &\subseteq\Core_E\left(\bigcup_{i\in I}B_i\right).
\end{aligned}
\]
Der Durchschnitt der leeren Familie wird hier relativ zu $A$
als $A$ verstanden. Die letzten beiden Inklusionen sind im
allgemeinen keine Gleichheiten.
\end{theorem}

\begin{theorem}[Mengenoperationen mit gesättigten Mengen]
Sei $E\in\EqSet(A)$.
Beliebige Vereinigungen und beliebige, relativ zu $A$ gebildete
Durchschnitte gesättigter Mengen sind gesättigt.
Auch Komplemente in $A$ und Differenzen gesättigter Mengen
sind gesättigt. Für $B\subseteq A$ und gesättigtes $D\subseteq A$
gilt zusätzlich
\[
 \Sat_E(B\cap D)=\Sat_E(B)\cap D,\qquad
 \Sat_E(B\setminus D)=\Sat_E(B)\setminus D.
\]
\end{theorem}

\section{Gesättigte Mengen und Teilfamilien der Quotientenmenge}

\begin{theorem}[Eindeutige Darstellung durch Klassen]
\label{b11guide:thm:subfamilies}
Für $E\in\EqSet(A)$ und $\mathcal D\subseteq A/E$ gilt
\[
 \{C\in A/E\mid C\subseteq\bigcup\mathcal D\}=\mathcal D.
\]
Für jede gesättigte Menge $B\subseteq A$ gilt
\[
 \bigcup\{C\in A/E\mid C\subseteq B\}=B.
\]
Insbesondere ist die Darstellung von $B$ als Vereinigung
einer Teilfamilie von $A/E$ eindeutig. Für Teilfamilien
$\mathcal C,\mathcal D\subseteq A/E$ gilt
\[
\begin{aligned}
 \bigcup\mathcal C\subseteq\bigcup\mathcal D
     &\ \Longleftrightarrow\ \mathcal C\subseteq\mathcal D,\\
 (\bigcup\mathcal C)\cap(\bigcup\mathcal D)
     &=\bigcup(\mathcal C\cap\mathcal D),\\
 A\setminus\bigcup\mathcal D
     &=\bigcup((A/E)\setminus\mathcal D).
\end{aligned}
\]
\end{theorem}

